\section{Description}
\label{sec:architecture}

%# Descrizione
%> Copiando da un paper di Cucinotta, qui descriverei solo il modello di cluster 
%    e l'API di ogni nodo sia per i client che per nodi.
%
%**No coordinazione.**
%Altre varie assunzioni:
%- Modello di cluster (possono esplodere per tempo finito).
%- Modello di banda.
%- No Consistenza: si fa con un layer aggiuntivo
%- API client + tabelle senza nome

The database is organized into \emph{tables}, each identified by a version~7 UUID~\cite{rfc9562}.
Every table contains key-value pairs, where both keys and values are byte arrays.

A Kitsurai system operates on a \emph{cluster} composed of a static set of nodes, which may occasionally be unavailable.
A \emph{node} corresponds to an instance of the program running on a physical or virtual machine.
Each node can be tagged with a string identifying its \emph{Availability Zone} (see Section~\ref{subsec:availability-zones}).
Every node maintains a copy of the metadata for all tables.

\subsection{Client API}
\label{subsec:interface}

Four main operations are supported:
\begin{itemize}
    \item \texttt{create-table}: creates a new table with the specified parameters.
    \item \texttt{delete-table}: deletes an existing table.
    \item \texttt{get}: retrieves the value associated with a key in a table.
    \item \texttt{set}: sets the value associated with a key in a table.
\end{itemize}

Any of these requests can be sent to a node chosen by the client, which acts as a \emph{router} to forward the request to the nodes responsible for the specified table.

\subsection{Tables}
\label{subsec:tables}

Each table has four parameters specified by the user at creation time:
\begin{itemize}
    \item \texttt{bandwidth} or \texttt{b}: the total number of requests (\texttt{get} and \texttt{set}) that the system must guarantee for the table, expressed in requests per second.
    \item \texttt{n}: the replication factor for each key-value pair.
    \item \texttt{read\_concern} or \texttt{r}: the minimum number of nodes that must successfully respond to a read request.
    \item \texttt{write\_concern} or \texttt{w}: the minimum number of nodes that must successfully respond to a write request.
\end{itemize}

And another four parameters computed automatically:
\begin{itemize}
    \item \texttt{id}: a UUIDv7 identifying the table, generated at creation.
    \item $\Sigma$: the effective bandwidth of the table considering the replication effect, computed as $\displaystyle \Sigma = b \cdot n$.
    \item \texttt{responsible\_nodes}: the list of nodes on which the table's key-value pairs are distributed, chosen to guarantee replication and data availability.
    \item $b_i$: the bandwidth allocated by node $i$ for the table, such that $\displaystyle \sum b_i = \Sigma$.
\end{itemize}

Each table is managed by a set of nodes, called \emph{responsible nodes}, chosen to guarantee replication and data availability.
On each node, a table can be in one of three states: \emph{prepared}, \emph{created}, or \emph{deleted}.
Tables can only progress in state in the order: \emph{prepared} $\rightarrow$ \emph{created} $\rightarrow$ \emph{deleted}.
A table is \emph{prepared} when the requested bandwidth has been allocated, but it has not yet been created.
A table is considered \emph{created} when all responsible nodes consider it so; it is considered \emph{deleted} when at least one node marks it as such.
The state of a table is synchronized among nodes through the \emph{Gossip} protocol (see Section~\ref{subsec:gossip}).

\subsection{Consistency model}
\label{subsec:consistency-model}

The system does not guarantee any form of consistency: there is no synchronization mechanism among nodes.
If a write does not reach a node because it is unavailable, that node will continue to return the previous value until it receives a new write.

Consistency guarantees can be strengthened by using the condition $r + w > n$.
In that case, at least one node involved in the read will have received the latest write, thus ensuring more consistent reads.
However, it is the client's responsibility to determine which of the returned values is the most recent, since the system does not provide timestamps or version identifiers.

\subsection{Availability Zones}
\label{subsec:availability-zones}

To ensure database availability, each node can be tagged with a string identifying its \emph{Availability Zone} (\texttt{AZ}).
Two nodes belonging to the same \texttt{AZ} cannot be chosen to replicate the same key.
To ensure that keys are replicated in a way that keeps them available even in case of an entire \texttt{AZ} failure, we modify the key distribution algorithm:
\begin{itemize}
    \item during table creation, the router must not allocate more than \texttt{b} bandwidth from nodes in the same \texttt{AZ};
    \item nodes must be assigned to the ring in order of \texttt{AZ}, so that an interval of length \texttt{n} in the ring contains nodes belonging to distinct \texttt{AZ}s.
\end{itemize}

\subsection{Throughput throttling}
\label{subsec:throughput-throttling}

In a multi-tenant system such as Kitsurai, it is necessary to prevent a single client from monopolizing resources and compromising the performance of others.
For this reason, Kitsurai implements a throughput throttling mechanism that guarantees each client the bandwidth requested at table creation, even in the presence of high and concurrent loads.

Each node in the cluster maintains, for each table, a variable \texttt{sched\_at} that indicates when the node will be available again to execute new operations on that table.
When a client sends a read or write request, the node waits until \texttt{sched\_at} before executing the operation, then increments \texttt{sched\_at} by an interval equal to $\displaystyle \frac{1}{b}$, where $b$ represents the requested bandwidth for the table.
If \texttt{sched\_at} is already in the past, the system could execute uncontrolled operations until \texttt{sched\_at} catches up with the current time; to avoid this, if \texttt{sched\_at} is already in the past, it is updated to 100 ms before the current time, thus allowing a short controlled burst of operations.

\subsection{Bandwidth allocation}
\label{subsec:bandwidth-allocation}

For bandwidth allocation, Kitsurai adopts a very simple model: it assumes that the maximum number of operations per second a machine can execute is fixed and independent of other factors such as the size of ongoing reads and writes.

For this reason, each node only needs to know its maximum bandwidth---measured beforehand using standard profiling and sizing operations---to determine whether it is able to handle the load requested by a table.
For this purpose, SSD disks must be used, as they provide constant and predictable performance, unlike HDDs, which are subject to unpredictable slowdowns due to seek operations and access latency.
